\documentclass[11pt]{article}
\usepackage[margin=1in]{geometry}
\usepackage{booktabs}
\usepackage{amsmath}
\usepackage{amssymb}
\usepackage{hyperref}
\usepackage{cite}

\title{Sex-Specific Resilience of Neocortex to Food Restriction}
\author{}
\date{}

\begin{document}
\maketitle

\begin{abstract}
Mammals exhibit sex-specific adaptations to reduce energy usage in times of food scarcity. Peripheral adaptations are well characterized, but much less is known about how the energy-expensive brain adapts and how such adaptations differ across sexes. We examined how food restriction (FR) leading to 15\% body-weight loss over 2--3 weeks impacts energy usage and function in primary visual cortex (V1) of adult male and female mice. Serum leptin was significantly decreased in males ($-$72\%, CTR vs FR: 6.45 ng/mL [95\% CI: 4.47--8.43] vs 1.84 ng/mL [95\% CI: 0.90--2.79]; $t = 4.58$, $df = 66$, $p < 0.0001$, $n = 17$ CTR males, 19 FR males) but not in females ($-$28\%, 3.94 vs 2.83 ng/mL; $t = 1.00$, $df = 66$, $p = 0.32$, 23 CTR females, 11 FR females). AMPK Thr172 phosphorylation in V1 was significantly elevated 2.9-fold in males ($t = 2.28$, $df = 39$, $p = 0.022$) but only 1.4-fold (n.s.) in females ($t = 0.64$, $df = 39$, $p = 0.11$). PPAR$\alpha$ DNA-binding activity in V1 was significantly regulated in males ($t = 4.81$, $df = 30$, $p = 0.013$) but not in females ($t = 0.0016$, $df = 30$, $p = 0.99$). RNA-seq of V1 identified 657 and 585 targets significantly differentially expressed in males and females respectively (p-adj $<$ 0.1), with 125 shared; GSEA identified 34 and 14 significantly enriched gene sets, with 13 shared. In vivo two-photon ATP imaging showed a 24\% decrease in FRET decay half-time in males ($t = 2.87$, $df = 37$, $p = 0.0067$) and 12\% non-significant decrease in females ($t = 1.36$, $df = 37$, $p = 0.18$). Orientation selectivity index decreased 27\% in males (CTR 0.58, FR 0.42; $t = 3.75$, $df = 27$, $p = 0.0009$) and 13\% non-significantly in females (CTR 0.60, FR 0.52; $t = 1.90$, $df = 27$, $p = 0.08$). Bayes factor analysis confirmed that food restriction effects were more robust in males.
\end{abstract}

\section{Introduction}
Mammals reduce energy usage during food scarcity, and peripheral adaptations are sex-specific: females readily lose muscle and bone mass but retain fat mass and body weight more than males \cite{lamont2008gender,fontana2009aging,manzanares2019cognitive,palmer2017sexual,mitchell2016effects,shi2020sex}, and preferentially suppress reproductive function (ovarian/uterine mass, cycling) \cite{wade1996metabolic,nelson1997reproduction,caligioni2009assessing}. In contrast, whether and how the brain adapts its function and energy usage in a sex-specific manner remains largely unknown. The brain consumes substantial energy \cite{attwell2001energy,harris2012synaptic,niven2016neuronal}, and the cerebral cortex represents over half of this \cite{harris2012synaptic}.

Previously we showed that food restriction in male mice reduced V1 energy usage during visual processing and decreased orientation selectivity, and that this required reduced serum leptin \cite{padamsey2022neocortex}. It is unclear whether females exhibit an analogous brain-level adaptation. Understanding sex-specific effects of food restriction is also critical for behavioral neuroscience, where food restriction is commonly used to motivate animals \cite{hori1990comparative,toth2000animal,guo2014procedures}.

Here we examined the molecular and functional consequences of a 2--3 week FR protocol (15\% body-weight loss) on V1 of adult male and female C57BL/6 mice, combining serum assays, western blotting, RNA-seq, in vivo FRET-based ATP imaging in Thy1-ATeam1.03$^{\mathrm{YEMK}}$ mice \cite{imamura2009visualization,yoshida2019development}, and two-photon GCaMP6s calcium imaging during drifting gratings.

\section{Methods}
\subsection{Animals and food restriction}
Adult C57BL/6 mice 7--9 weeks old were given ad libitum (CTR) access to food or food-restricted (FR) to 85\% of free-feeding body weight for 2--3 weeks. Prior to recordings and tissue collection, all animals were given ad libitum access to food until sated. To achieve similar magnitude (15\%) and rate of weight loss in both sexes, FR females received 39\% of ad libitum intake on average whereas FR males received 27\%, a 40\% greater restriction in females.

\subsection{Serum leptin and western blotting}
Serum leptin was measured by ELISA. V1 tissue was probed for AMPK Thr172 phosphorylation normalized by total AMPK, and PPAR$\alpha$ activity was assessed as DNA-binding normalized to protein level.

\subsection{RNA sequencing}
RNA-seq was performed on V1 tissue ($n = 4$ per CTR/FR per sex). Differential expression analysis used DESeq2 \cite{love2014moderated} with p-adj $< 0.1$; Gene Set Enrichment Analysis used Hallmark Gene Sets \cite{subramanian2005gene,liberzon2015molecular}.

\subsection{ATP imaging}
Two-photon FRET imaging of Thy1-ATeam1.03$^{\mathrm{YEMK}}$ \cite{imamura2009visualization,yoshida2019development} in awake, head-fixed mice during presentation of natural outdoor scenes was used. ATP synthesis inhibitors applied focally over V1 unmasked ATP usage as decay of the FRET signal. Half-time of FRET decay was quantified.

\subsection{Two-photon calcium imaging}
V1 L2/3 GCaMP6s \cite{chen2013ultrasensitive}-labeled neurons were imaged during presentation of drifting gratings. Orientation selectivity index (OSI) was computed. Male data were compared to previously published datasets obtained under similar conditions \cite{padamsey2022neocortex}.

\subsection{Statistics}
Two-way ANOVA was used for group comparisons, with 95\% CIs. Bayes factor BF$_{10}$ analysis \cite{kass1995bayes,morey2015baws} quantified relative evidence for the alternative vs null hypothesis.

\section{Results}
\subsection{Weight loss, food intake, and serum leptin}
Ad libitum daily food intake was comparable between sexes despite females weighing $\approx$80\% of age-matched males (male vs female: 27.82 g [95\% CI: 26.84--28.79] vs 23.21 g [95\% CI: 22.35--24.07]; $t = 7.58$, $df = 66$, $p < 0.001$, $n = 8$ male, 7 female). FR required 40\% greater caloric restriction in females (39\% of ad libitum) than males (27\% of ad libitum) to achieve 15\% body-weight loss. Daily food intake two-way ANOVA: CTR male vs FR male, $t = 4.81$, $df = 25$, $p < 0.0001$; CTR male vs CTR female $t = 0.013$, $df = 25$, $p = 0.99$; CTR female vs FR female $t = 6.59$, $df = 25$, $p < 0.0001$; FR male vs FR female $t = 1.93$, $df = 25$, $p = 0.07$. Animal weight two-way ANOVA: CTR male vs FR male $t = 11.36$, $df = 25$, $p < 0.0001$; CTR female vs FR female $t = 8.36$, $df = 25$, $p < 0.0001$.

Serum leptin decreased $-$72\% in males (CTR 6.45 ng/mL [95\% CI: 4.47--8.43] vs FR 1.84 ng/mL [95\% CI: 0.90--2.79]; $t = 4.58$, $df = 66$, $p < 0.0001$) and only $-$28\% (n.s.) in females (CTR 3.94 [95\% CI: 2.53--5.35] vs FR 2.83 [95\% CI: 1.20--4.47]; $t = 1.00$, $df = 66$, $p = 0.32$).

\begin{table}[h]
\centering
\caption{Body weight, food intake, and serum leptin (two-way ANOVA).}
\begin{tabular}{llll}
\toprule
Measure & Comparison & $t$, $df$, $p$ & $n$ \\
\midrule
Daily food intake & CTR male vs FR male & $t = 4.81$, $df = 25$, $p < 0.0001$ & 8/8/7/7 \\
Daily food intake & CTR female vs FR female & $t = 6.59$, $df = 25$, $p < 0.0001$ & \\
Daily food intake & FR male vs FR female & $t = 1.93$, $df = 25$, $p = 0.07$ & \\
Body weight & CTR male vs FR male & $t = 11.36$, $df = 25$, $p < 0.0001$ & 17/19/23/11 \\
Body weight & CTR female vs FR female & $t = 8.36$, $df = 25$, $p < 0.0001$ & \\
Serum leptin (ng/mL) & CTR male (6.45) vs FR male (1.84) & $t = 4.58$, $df = 66$, $p < 0.0001$ & 17/19 \\
Serum leptin (ng/mL) & CTR female (3.94) vs FR female (2.83) & $t = 1.00$, $df = 66$, $p = 0.32$ & 23/11 \\
\bottomrule
\end{tabular}
\end{table}

\subsection{V1 AMPK and PPAR$\alpha$ are regulated by FR in males but not females}
AMPK Thr172 phosphorylation was elevated 2.9-fold in males (CTR 13.19 AU/$\mu$g [95\% CI: 7.68--18.71] vs FR 38.21 AU/$\mu$g [95\% CI: 16.50--59.93]; $t = 2.28$, $df = 39$, $p = 0.022$, $n = 11$ CTR, 15 FR males), and only 1.4-fold in females (CTR 20.94 AU/$\mu$g [95\% CI: 9.42--32.46] vs FR 29.52 AU/$\mu$g [95\% CI: 8.52--50.53]; $t = 0.64$, $df = 39$, $p = 0.11$, $n = 11$ CTR, 6 FR females). PPAR$\alpha$ activity was significantly regulated in males ($t = 4.81$, $df = 30$, $p = 0.013$, $n = 5$ CTR, 11 FR) but not females ($t = 0.0016$, $df = 30$, $p = 0.99$, $n = 9$ CTR, 9 FR).

\subsection{RNA-seq reveals male-specific regulation of energy pathways}
DESeq2 identified 657 targets in males and 585 in females at p-adj $< 0.1$ (125 shared). GSEA (Hallmark) revealed 34 (males) vs 14 (females) significantly regulated gene sets (p-adj $< 0.05$; 13 shared). In males, gene sets associated with oxidative phosphorylation ($t = 0.48$, $p < 0.0001$), mTORC1 signaling ($t = 0.46$, $p < 0.0001$), fatty acid metabolism ($t = 0.43$, $p < 0.0001$), and PI3K/AKT/mTOR signaling ($t = 0.34$, $p = 0.027$) were significantly regulated; none of these were significantly regulated in females (oxidative phosphorylation $t = -0.17$, $p = 0.96$; mTORC1 $t = -0.18$, $p = 0.98$; fatty acid $t = 0.2$, $p = 0.88$; PI3K/AKT/mTOR $t = 0.17$, $p = 0.98$). Estrogen-regulating pathways were also regulated in males but not females.

\begin{table}[h]
\centering
\caption{RNA-seq and GSEA summary (V1 tissue).}
\begin{tabular}{lll}
\toprule
Measure & Males & Females \\
\midrule
Differentially expressed genes (p-adj $<$ 0.1) & 657 & 585 \\
Shared differentially expressed genes & 125 & 125 \\
GSEA significant gene sets (p-adj $<$ 0.05) & 34 & 14 \\
Shared significant gene sets & 13 & 13 \\
Oxidative phosphorylation & $t = 0.48$, $p < 0.0001$ & $t = -0.17$, $p = 0.96$ \\
mTORC1 signaling & $t = 0.46$, $p < 0.0001$ & $t = -0.18$, $p = 0.98$ \\
Fatty acid metabolism & $t = 0.43$, $p < 0.0001$ & $t = 0.2$, $p = 0.88$ \\
PI3K/AKT/mTOR signaling & $t = 0.34$, $p = 0.027$ & $t = 0.17$, $p = 0.98$ \\
\bottomrule
\end{tabular}
\end{table}

\subsection{ATP usage in V1 is reduced by FR in males but not females}
Half-time of FRET decay during presentation of natural outdoor scenes was slowed by 24\% in males (CTR 8.48 min [95\% CI: 7.59--9.37] vs FR 10.51 min [95\% CI: 9.08--11.93]; $t = 2.87$, $df = 37$, $p = 0.0067$, $n = 11$ CTR, 10 FR), indicating reduced ATP usage. Females showed a modest 12\% non-significant change (CTR 8.66 [95\% CI: 7.49--9.84] vs FR 9.66 [95\% CI: 8.92--10.39]; $t = 1.36$, $df = 37$, $p = 0.18$, $n = 12$ CTR, 8 FR). In darkness, ATP usage was unaffected for either sex (males $t = 0.14$, $df = 24$, $p = 0.89$, $n = 8/6$; females $t = 0.64$, $df = 37$, $p = 0.52$, $n = 8/8$).

\subsection{Orientation and direction selectivity are reduced in males only}
Orientation selectivity index (OSI) decreased 27\% in males (CTR 0.58 [95\% CI: 0.51--0.65] vs FR 0.42 [95\% CI: 0.36--0.49]; $t = 3.75$, $df = 27$, $p = 0.0009$, $n = 8$ CTR, 8 FR), and 13\% non-significantly in females (CTR 0.60 [95\% CI: 0.52--0.67] vs FR 0.52 [95\% CI: 0.45--0.58]; $t = 1.90$, $df = 27$, $p = 0.08$, $n = 7$ CTR, 9 FR). Direction selectivity index was also significant in males (CTR vs FR: $t = 2.16$, $p = 0.023$, $df = 12$, $n = 6/8$) but not females ($t = 1.49$, $p = 0.16$, $df = 14$, $n = 7/9$). Pupil diameter did not differ across sex and diet ($t = 0.82$, $df = 18$, $p = 0.42$).

\begin{table}[h]
\centering
\caption{V1 ATP usage and visual coding (two-way ANOVA).}
\begin{tabular}{llll}
\toprule
Measure & CTR vs FR (males) & CTR vs FR (females) & Notes \\
\midrule
FRET decay half-time (natural scenes) & 8.48 vs 10.51 min; $t=2.87$, $p=0.0067$ & 8.66 vs 9.66 min; $t=1.36$, $p=0.18$ & 24\% vs 12\% \\
FRET decay half-time (darkness) & $t=0.14$, $p=0.89$ & $t=0.64$, $p=0.52$ & n.s. \\
Orientation selectivity index & 0.58 vs 0.42; $t=3.75$, $p=0.0009$ & 0.60 vs 0.52; $t=1.90$, $p=0.08$ & 27\% vs 13\% \\
Direction selectivity index & $t=2.16$, $p=0.023$ & $t=1.49$, $p=0.16$ & \\
\bottomrule
\end{tabular}
\end{table}

\subsection{Bayes factor analysis}
Across the 10 parameters of energy usage and visual coding tested, BF$_{10}$ for males was consistently several-fold greater than for females, indicating that FR produced more robust effects in males. In females, 3 of 10 parameters had BF$_{10} < 0.33$, consistent with a small residual impact below statistical detection.

\section{Discussion}
Consistent with peripheral studies, we find that V1 energy usage and visual coding are largely resistant to FR in female mice, while they are robustly reduced in males. Reduced serum leptin is central to the male response \cite{padamsey2022neocortex}, and its attenuated decrease in females likely explains the preservation of cortical function. Estrogen and progesterone may additionally suppress food intake and promote energy expenditure \cite{brown2015estrogen}, augmenting leptin signalling via STAT3 phosphorylation \cite{gao2007estrogens} and upregulating mitochondrial function \cite{klinge2008mitochondrial,ventura2018sex}. In our RNA-seq, estrogen signalling pathways were regulated in males but not females, suggesting overlapping effects between estrogen and FR that could diminish further female responses.

AMPK and mTOR are canonical cellular energy regulators \cite{mihaylova2011ampk,saxton2017mtor}, and their sex- and tissue-dependent regulation by FR is established in peripheral tissues \cite{ventura2018sex,goldspink2019exercise,kenny2013endurance}. Our study extends this sex-specific regulation to the neocortex. Because mTOR regulates neuronal excitability, neurite outgrowth, ion channel expression, and synaptic plasticity \cite{saxton2017mtor}, its FR-driven inhibition in males could underlie reduced ATP usage and loss of orientation selectivity.

Limitations include the moderate 2--3 week FR protocol, single cortical region, and lack of statistical power to directly compare effect sizes between sexes (which would require much larger groups). Bayes factor analysis supports stronger male effects but cannot conclusively rule out a small effect in females. Sex differences in FR effects on adult cognition have been reported \cite{hallet2017juvenile,hallet2019juvenile}; our findings are relevant both to studies of FR as a motivational tool and to dietary interventions in neurological disease \cite{mattson2017calorie,gano2014ketogenic}. Including both sexes in biomedical research is essential to ensure equitable treatment development.

\bibliographystyle{plain}
\bibliography{references}

\end{document}
